\section{Experience Library Examples}


\subsection{Mathematical Reasoning}

\begin{verbatim}
[G0]. When solving geometry problems with intersections,
      validate solutions lie within bounded regions, not
      extensions, to avoid extraneous answers.

[G1]. For expected extreme statistics in combinatorial
      problems, use direct enumeration for small sizes.

[G10]. When using mathematical invariants to prove
       impossibility, always validate against known
       achievable states or small cases.

[G21]. For complex polynomials with real parameters,
       separate real/imaginary parts to find when
       real roots exist.
\end{verbatim}

\subsection{Code Generation}

\begin{verbatim}
[C0]. Before implementing complex algorithms, write
      unit tests for edge cases (empty input, single
      element, maximum size).

[C5]. When optimizing nested loops, consider whether
      inner loop can be vectorized or replaced with
      hash map lookup.

[C12]. For recursive functions, always define base
       case first and verify termination condition.
\end{verbatim}

\subsection{Creative Writing}

\begin{verbatim}
[W0]. When developing character arcs, establish clear
      motivation in first act to justify later
      decisions.

[W3]. In descriptive passages, prefer specific sensory
      details over generic adjectives.

[W7]. For dialogue, ensure each character has distinct
      speech patterns reflecting background/personality.
\end{verbatim}

